\section{Conclusions}

We compared three VQA architectures across two benchmarks, covering a decade of progress from CNN+LSTM fusion to generative pretrained models. On VQA v2 the progression is unambiguous: 56.88\% $\rightarrow$ 70.30\% $\rightarrow$ 75.95\%. On GQA the ranking reverses, with LXMERT (59.29\%) outperforming BLIP (47.11\%) by a wide margin.

Three conclusions emerge. First, cross-modal attention (the transition from CNN-RNN fusion to LXMERT) accounts for the largest single performance jump on VQA v2, more than the subsequent shift to generation, suggesting that fine-grained alignment between image regions and question tokens is the central challenge in VQA. Second, generative pretraining does not automatically transfer across datasets: without GQA-specific fine-tuning, BLIP's open-vocabulary answers no longer match GQA's expected answer distribution, and BLIP falls below even the Vanilla baseline. Third, at least Vanilla VQA and LXMERT appear to work mainly by picking the answer that matched similar patterns during training, rather than truly reasoning about the image and question, and BLIP's own collapse on GQA suggests it relies on the same kind of pattern matching to some extent.
